\section{Design Tradeoffs}

Ghostext prioritizes deterministic correctness and fail-closed behavior before maximizing payload capacity. This section explains the major tradeoffs and how each one should be interpreted in practice.

\subsection{Capacity vs. Synchronization Safety}

Wider candidate sets can increase available interval partitioning and improve payload efficiency. The same change can also increase exposure to tokenization instability and path divergence. Ghostext therefore keeps candidate-policy knobs explicit (top-$p$, max-$k$, min-entropy threshold) so operators can tune capacity only with visible safety consequences.

\subsection{Fluency vs. Invertibility}

Unconstrained generation can improve surface fluency but can break deterministic replay. Ghostext separates payload-carrying prefix from optional non-payload tail to keep invertibility-critical logic isolated. This separation does not remove all detectability concerns, but it prevents fluency tuning from silently weakening decode correctness.

\subsection{Truncation Efficiency vs. Distribution Fidelity}

Truncation reduces runtime and narrows search but introduces a policy-dependent statistical gap from base model output. In Ghostext this effect is tracked by $\alpha_t$, while quantization and stability filtering add the other terms in the Section~5 bound:
\begin{equation}
\mathrm{TV}(Q'_t,P_t)\le\alpha_t + \frac{N_t}{2F_{\mathrm{tot}}} + \beta_t.
\end{equation}

The practical implication is that faster and narrower policies should be treated as explicit statistical tradeoffs, not as free optimizations.

\subsection{Retry Utility vs. Traffic Leakage}

Retry logic can recover from low-entropy or unstable-tokenization regions and improve decode success in difficult contexts. The same mechanism can alter observable latency and output-length distributions. Ghostext keeps retries bounded and classifies retry-triggered failures so this tradeoff remains measurable.

\subsection{Strict Mismatch Rejection vs. Backward Compatibility}

Configuration fingerprint checks aggressively reject runtime drift. This prevents silent corruption and ambiguous partial decode results. The cost is operational rigidity: upgrades that change backend metadata or tokenizer identity require coordinated rollout. For security-sensitive use, Ghostext chooses strict rejection over permissive compatibility.

\subsection{Deterministic Policies vs. Adaptive Auto-Tuning}

Adaptive online tuning might improve local efficiency, but it complicates reproducibility and can obscure causal attribution in security evaluation. Current design favors deterministic policies with audit-friendly logs. Future adaptive extensions should only be introduced with explicit replayability constraints and side-channel monitoring.

\subsection{Granular Failure Taxonomy vs. Simplicity}

A single generic failure code is simpler to expose, but hides important operational and security differences. Ghostext adopts a richer failure taxonomy to support debugging, deployment tuning, and honest reporting. The overhead is additional implementation complexity and more detailed evaluation requirements.

\subsection{Two-Stage Packet Embedding vs. Monolithic Embedding}

A monolithic single-segment design is conceptually simpler, but delays detection of protocol mismatches. Two-stage embedding adds protocol structure and some overhead, yet enables early authenticated checks and clearer diagnosis. For the current engineering goals, this extra structure is a deliberate choice.

\subsection{Constrained Tail Reuse vs. Separate Tail Sampler}

Reusing the constrained candidate mechanism for natural tails keeps code paths unified and predictable. A separate unconstrained tail sampler may produce different style characteristics but would add another source of distribution and reproducibility variance. Current implementation favors one-path simplicity while treating tail policy as an explicit ablation axis.

\subsection{Current Position}

None of these tradeoffs is universally optimal. Ghostext makes conservative choices that favor deterministic replay, explicit checks, and measurable behavior. Post-pause experiments should decide which knobs can be relaxed without unacceptable recoverability, side-channel, or detector-facing costs.
